% Solutions of 30-project-text-search-engine.
\section*{Chapter \ref{ch:b1:project-text-search-engine} --- Project: A Search Engine for a Folder of Texts}

\begin{solution}{exo:b1:project-text-search-engine:1}
\emph{it}, \emph{s}, \emph{nal}, \emph{s}, \emph{e}, \emph{mail}, \emph{v}. The digits and punctuation only separate.
The letter Ü is the two bytes \code{0xC3 0x9C}, neither of them an ASCII letter, so it separates too. A reader would not
call \emph{s}, \emph{nal}, \emph{e} or \emph{v} words. The query \emph{Ünal} goes through the same \omterm{def:b1:project-text-search-engine:tokenizer}{tokenizer} and becomes
\emph{nal}, which is in the index for this line. So it is found, together with every other text in which a non-ASCII
letter is followed by \emph{nal}: the rule is consistent, not exact.
\end{solution}

\begin{solution}{exo:b1:project-text-search-engine:2}
\emph{the}: $(0, 2), (1, 2)$. \emph{cat}: $(0, 1), (1, 1), (2, 1)$. \emph{hat}: $(1, 1), (2, 1)$. \emph{a}: $(2, 3)$.
The capitals of document 1 are lowercased, and the comma and semicolon of document 2 separate. There are 9 distinct
words (\emph{the, cat, sat, on, mat, and, hat, a, bat}) and $5 + 4 + 4 = 13$ postings, one per distinct word of each
document.
\end{solution}

\begin{solution}{exo:b1:project-text-search-engine:3}
The steps, written (document of the result list, document of the second list): $(2, 1)$ advances the second list;
$(2, 5)$ the first; $(5, 5)$ keeps 5; $(9, 6)$ the second; $(9, 14)$ the first; $(14, 14)$ keeps 14; $(20, 31)$ the
first; $(31, 31)$ keeps 31. The result is $[5, 14, 31]$, after 8 runs of the body. The first list is then exhausted,
and since every later document of the second list (40 and 41) is greater than all of the first, none of them can be
in both.
\end{solution}

\begin{solution}{exo:b1:project-text-search-engine:4}
Writing hits by their scores: $[4]$; $[4, 8]$, since 8 is better than its parent 4 and stays below it; $[1, 8, 4]$, the
1 rising to the root; then the heap is full. 6 beats the root 1 and replaces it, and sinks below its worse child 4:
$[4, 8, 6]$. 9 beats 4, replaces it and sinks below 6: $[6, 8, 9]$. 3 does not beat 6 and is dropped. 7 beats 6 and
stays at the root, both children being better: $[7, 8, 9]$. \code{topk_sorted} returns 9 (document 4), 8 (document 1)
and 7 (document 6).
\end{solution}

\begin{solution}{exo:b1:project-text-search-engine:5}
A word is counted only when a separator ends it. In the third text the last word, \emph{fish}, is followed by no byte
at all, so it is never counted. The function fails on every text whose last byte is a letter, which is a file without
a final newline ending on a word, or any string passed without one. The correction counts the word still open at the
end: \code{return n + in_word;}. \code{next_word} has no such case, because its loop over letters stops at the end of
the text as well as at a separator, and returns the word it has copied.
\end{solution}

\begin{solution}{exo:b1:project-text-search-engine:6}
Every occurrence of \emph{cat} adds a posting, so document 0, which has it twice, gets two postings $(0, 1)$. The
length of the list, taken as the \omterm{def:b1:project-text-search-engine:tfidf}{document frequency}, is then 3, more than the 2 documents. The weight $\ln(2/3)$ is
negative, so a document would lose score for containing the word. In the engine, line 62 of
\cref{lst:b1:project-text-search-engine:index} looks at the last posting and increments its count when it is already
for the current document. With the error, the engine's document frequencies would count occurrences, its scores would
differ from those of the naive scan, which counts documents, and the \omterm{def:b1:project-text-search-engine:end-to-end}{end-to-end test} would fail on the first query
of a word used twice in a document, that is, on almost every query.
\end{solution}

\begin{solution}{exo:b1:project-text-search-engine:7}
An \omterm{def:b1:pointers-and-arrays:array}{array} of scores, one per document, receives the weight of each posting of each term; a document is a match if it
received anything. The matches then go through the heap.

\omcode{code/bachelor-1/30-project-text-search-engine/exercises/or_query.c}{7}{34}{OR queries, by adding the weights of
every posting}

The test \code{tests/exercises_test.c} compares the number of matches with a brute-force count on six small documents,
for every pair of six words, one of which occurs nowhere. The cost is the total length of the \omterm{def:b1:project-text-search-engine:inverted-index}{posting lists},
$\ell_1 + \dots + \ell_k$, plus $\Oh{N}$ to clear and scan the \omterm{def:b1:pointers-and-arrays:array}{arrays}. There is no shortcut by the shortest list as for
AND: every posting of every list contributes. The result can be large too, since an OR query with one frequent word
matches almost every document.
\end{solution}

\begin{solution}{exo:b1:project-text-search-engine:8}
The function takes the shortest list and keeps each of its documents only if a binary search finds it in every other
list; it counts and allocates nothing.

\omcode{code/bachelor-1/30-project-text-search-engine/exercises/count.c}{7}{35}{Counting the documents that contain
every term}

In \code{main}, when \code{argv[2]} is \texttt{-{}-{}count}, the terms are split from \code{argv + 3} instead, and the program
prints \code{count_and(&ix, terms, n)} and returns 0 if that number is positive, 1 otherwise. The test compares
\code{count_and} with the brute-force count of the previous exercise.
\end{solution}

\begin{solution}{exo:b1:project-text-search-engine:9}
\begin{enumerate}
\item The running result starts as the shortest list and only loses documents, so it has at most $\ell_1$ elements.
In the merge with list $i$, every run of the body advances the index of the result, or that of list $i$, or both. The
first index takes at most $\ell_1$ values and the second at most $\ell_i$, so the body runs at most $\ell_1 + \ell_i$
times. Summing over $i = 2, \dots, k$ gives $(k - 1)\ell_1 + \ell_2 + \dots + \ell_k$.
\item Each of the $\ell_1$ documents is searched for in $k - 1$ lists, and a binary search in a list of length $\ell_i
\le \ell_k$ makes at most $\lfloor \log_2 \ell_i \rfloor + 1 \le \log_2 \ell_k + 1$ comparisons. In all, $\ell_1 (k -
1)(\log_2 \ell_k + 1) = \Oh{\ell_1 \log \ell_k \cdot k}$.
\item \emph{xylograph} is in one document (42) and \emph{ba} in 400. The bound of the merge is $1 + 400 = 401$ runs; the
real merge runs 43 times, advancing through \emph{ba}'s documents 0 to 42, then stops when the result is exhausted.
The binary search makes at most $\lfloor \log_2 400 \rfloor + 1 = 9$ comparisons. When one list is much shorter than
the others, searching beats merging; when they have similar lengths, $\ell_1 \log \ell_k$ exceeds $\ell_1 + \ell_k$
and merging wins.
\end{enumerate}
\end{solution}

\begin{solution}{exo:b1:project-text-search-engine:10}
A positional \omterm{def:b1:project-text-search-engine:inverted-index}{posting list} holds, for a word, the pairs (document, position) of its occurrences, in increasing order.
The phrase starts at the places of $w_1$ that are followed by $w_i$ at $+ (i - 1)$ for every $i$. Each word filters the
places kept so far by one merge, comparing $(d, p + i - 1)$ with the places of $w_i$ in \omterm{def:b1:induction-and-termination:lex}{lexicographic order}.

\omcode{code/bachelor-1/30-project-text-search-engine/exercises/phrase.c}{40}{69}{Phrase queries, by merging lists of
(document, position) pairs}

Here the positional lists are built from the texts by \code{places_of}; in the engine they would be stored in each
posting, as an \omterm{def:b1:pointers-and-arrays:array}{array} of positions next to the count, which multiplies the size of the index by about the number of
occurrences per posting. The test finds \emph{red fish} in documents 0 and 3, and \emph{fish red} in none, although
both words are in documents 0 and 3: \emph{fish} is never just before \emph{red}.
\end{solution}

\begin{solution}{exo:b1:project-text-search-engine:11}
An entry is a \omterm{def:b1:pointers-and-arrays:pointer}{pointer} to its word, a \omterm{def:b1:project-text-search-engine:inverted-index}{posting list} (a \omterm{def:b1:pointers-and-arrays:pointer}{pointer} and two \code{int}s) and a \omterm{def:b1:pointers-and-arrays:pointer}{pointer} to the next entry:
$8 + 16 + 8 = 32$ bytes. A posting is two \code{int}s, 8 bytes. The \omterm{def:b1:pointers-and-arrays:array}{array} of names has a capacity $D$ that doubles from
64. Then
\[
  M = 8B + 32W + (L + W) + 8 \sum_t c_t + 8D + \sum_d (|\mathit{name}_d| + 1) .
\]
The table starts with 1024 \omterm{def:b1:hash-tables:hash-table}{buckets} and doubles when the number of words exceeds the number of \omterm{def:b1:hash-tables:hash-table}{buckets}: at the 1025th,
2049th and 4097th word, so $B = 8192$ for 5004 words. Likewise $D = 512$ for 400 documents. Then $8B = 65\,536$,
$32W = 160\,128$, $L + W = 35\,207$, $8 \sum_t c_t = 2\,192\,688$, $8D = 4096$ and the names take $400 \times 11 = 4400$
bytes: $M = 2\,462\,055$ bytes, the number printed. The posting \omterm{def:b1:pointers-and-arrays:array}{arrays} hold $274\,086$ places for $193\,455$ postings, so
$80\,631$ places, 29\,\% of them and 645\,048 bytes, are slack.
\end{solution}

\begin{solution}{exo:b1:project-text-search-engine:12}
The format: the number of documents and their names, then the number of words, then each word followed by the length
of its \omterm{def:b1:project-text-search-engine:inverted-index}{posting list} and its postings; every number is an \code{int}, every string its length then its bytes.

\omcode{code/bachelor-1/30-project-text-search-engine/exercises/persist.c}{32}{47}{Saving the index}

Loading rebuilds, for each document, a text made of its words, each repeated as often as its count, and indexes these
texts in document order: the lists, counts and numbers come out as before, through the code that is already tested.
The loader refuses a file that ends early (every \code{fread} is checked), a name or word longer than 4096 bytes, a
document number out of range, an empty list or a count that is not positive, and returns $-1$ without crashing. The test
in \code{tests/exercises_test.c} saves, loads, and compares the numbers of documents, words and postings and the
\omterm{def:b1:project-text-search-engine:inverted-index}{posting lists} of five words. The file is not portable because it stores \code{int}s and
\code{Posting}s as the machine lays them out in memory: a machine of the other \omterm{def:b1:c-strings-structs-files:byte-order}{byte order}
(\cref{def:b1:c-strings-structs-files:byte-order}) or with another size of \code{int} would read other numbers. A
portable format fixes the size and the \omterm{def:b1:c-strings-structs-files:byte-order}{byte order} of every number, as network formats do.
\end{solution}
